\documentclass[11pt,a4paper]{article}

% ── Packages ──────────────────────────────────────────────────────────────────
\usepackage[margin=2.5cm]{geometry}
\usepackage{amsmath,amssymb,amsthm}
\usepackage{graphicx}
\usepackage{booktabs}
\usepackage{hyperref}
\usepackage{xcolor}
\usepackage{microtype}
\usepackage{caption}
\usepackage{float}
\usepackage{enumitem}
\usepackage{setspace}
\usepackage{natbib}
\usepackage{array}
\usepackage{multirow}
\usepackage{longtable}
\usepackage{tabularx}
\usepackage{mdframed}
\usepackage{tcolorbox}
\usepackage{algorithm}
\usepackage{algpseudocode}

\hypersetup{
  colorlinks=true,
  linkcolor=blue!60!black,
  citecolor=green!40!black,
  urlcolor=blue!60!black,
}

\captionsetup{font=small, labelfont=bf}

% ── Color boxes ───────────────────────────────────────────────────────────────
\tcbuselibrary{skins,breakable}
\newtcolorbox{testedbox}{colback=green!6!white, colframe=green!50!black,
  fonttitle=\bfseries, title=Tested / Empirically Validated, breakable}
\newtcolorbox{plannedbox}{colback=orange!6!white, colframe=orange!60!black,
  fonttitle=\bfseries, title=Planned / Not Yet Run, breakable}
\newtcolorbox{scaffoldbox}{colback=blue!5!white, colframe=blue!40!black,
  fonttitle=\bfseries, title=Implemented as Scaffolding (Phase A), breakable}

% ── Theorems ──────────────────────────────────────────────────────────────────
\newtheorem{theorem}{Theorem}
\newtheorem{corollary}[theorem]{Corollary}
\newtheorem{definition}{Definition}

% ── Macros ────────────────────────────────────────────────────────────────────
\newcommand{\cliff}{\Delta_{\mathrm{cliff}}}
\newcommand{\wcliff}{\Delta_{W\text{-cliff}}}
\newcommand{\bcliff}{\Delta_{B\text{-cliff}}}
\newcommand{\kap}{\kappa}
\newcommand{\rhat}{\hat{r}}
\newcommand{\hhat}{\hat{h}}
\newcommand{\tauq}{\tau_q}
\newcommand{\CG}{\textsc{Cliffguard}}
\newcommand{\YES}{\textcolor{green!50!black}{\textbf{\checkmark}}}
\newcommand{\NO}{\textcolor{red!60!black}{\textbf{---}}}
\newcommand{\PART}{\textcolor{orange!70!black}{\textbf{$\sim$}}}

% ── Title ─────────────────────────────────────────────────────────────────────
\title{%
  \textbf{CLIFFGUARD: Research Proposal}\\[0.3em]
  {\large An Edge-Native, Quantization-Aware, Black-Box-Tolerant,
  RL-Adapted Defense System\\Against Prompt Injection at the Safety Cliff}\\[0.5em]
  {\normalsize \textit{Full framework blueprint, Phase A/B empirical results,
  and complete evaluation roadmap}}
}

\author{
  Trilochan Bhandari\\
  \small Department of Computer Science and Engineering\\
  \small Kathmandu University, Dhulikhel, Nepal\\
  \small \texttt{trilochan2625@student.ku.edu.np}
}

\date{May 2026}

% ─────────────────────────────────────────────────────────────────────────────
\begin{document}
\maketitle
\thispagestyle{plain}

\begin{abstract}
\CG{} is a stateless, online-RL-adapted, defense-in-depth system for prompt
injection on edge-deployed large language models (LLMs) operating across
heterogeneous post-training quantization (PTQ) regimes.
It is motivated by a body of evidence that PTQ degrades RLHF-installed safety
behavior non-linearly~\citep{egashira2024exploiting,mindthegap2025,hong2024quant}
and that the refusal direction in the residual stream can be probed as a
sensor~\citep{arditi2024refusal,zhao2025refusal}.

This document is the full research proposal for \CG{}.
It describes: (1) the complete system architecture and its eight named components;
(2) the formal theoretical foundations (decoupling theorem, CUSUM, LinUCB
formulation); (3) the threat model of nine adversaries; (4) the five
pre-registered hypotheses and the five-fold evaluation protocol; (5)
a precise account of what has been \textbf{implemented}, what has been
\textbf{tested with real data}, and what \textbf{remains planned};
(6) first empirical results from Folds~A and~B on Google Colab T4
(\texttt{Llama-3.2-3B-Instruct}, FP16 and NF4); (7) comparison to eleven
existing defense systems; and (8) timeline, resources, and
responsible-disclosure plan.

\textbf{Current empirical status (May 2026):}
Fold~A calibration and Fold~B cliff measurement are complete on one model family
and two quantization schemes.
Geometric cliff metric $\cliff = 0.167$ ($0.67\kap$, where $\kap=0.25$);
Wasserstein metric $\wcliff = 0.014$; behavioral proxy $\bcliff = 0.000$.
Hypothesis H1 (cliff existence) is not accepted for this configuration.
The predicted cliff zone (GGUF Q3\_K\_M and below) has not yet been evaluated.
Folds~C--E and Hypotheses H2--H5 are planned.
\end{abstract}

\setstretch{1.15}
\tableofcontents
\newpage

% ─────────────────────────────────────────────────────────────────────────────
\section{Motivation and Problem Statement}
\label{sec:motivation}
% ─────────────────────────────────────────────────────────────────────────────

\subsection{The Safety Cliff}

Modern instruction-tuned LLMs encode safety alignment via RLHF~\citep{ouyang2022training,
bai2022training}.
\citet{arditi2024refusal} showed this alignment is mediated by a single linear
direction $\rhat$ in the residual stream; ablating $\rhat$ raises attack
success rates dramatically.
\citet{zhao2025refusal} (NeurIPS 2025) extended this to show that
\emph{harmfulness} and \emph{refusal} are encoded in \textbf{separate} directions:
$\hhat$ at the user-instruction token, $\rhat$ at the post-instruction token.

Post-training quantization (PTQ)---reducing weights from FP16 to INT8, NF4, AWQ-INT4,
GGUF Q4/Q3/Q2---enables edge deployment but may destroy this safety subspace:

\begin{itemize}[leftmargin=2em, itemsep=3pt]
  \item \citet{hong2024quant} showed GPTQ-3-bit on Llama-2-13B-Chat drops toxicity
        safety by ${\sim}50$ points while MMLU falls only ${\sim}8$ points.
  \item \citet{egashira2024exploiting} (NeurIPS 2024) demonstrated the
        \emph{cliff-exploit attack}: adversaries craft FP16 model weights that
        appear safe, where malicious behavior activates \emph{only after} GGUF
        quantization, with ASR deltas of $\Delta = 88.7\%$ (insecure code),
        $\Delta = 85.0\%$ (content injection), and $\Delta = 30.1\%$ (refusal bypass).
  \item \citet{mindthegap2025} (ICLR 2025) showed GGUF k-quants' MSE objective
        does not preserve refusal-related activation geometry.
  \item \citet{wee2025caq} showed PTQ optimized for perplexity reconstruction
        systematically degrades behavioral alignment even when perplexity is preserved.
\end{itemize}

We define the \textbf{safety cliff} as the quantization level at which
refusal degradation crosses a pre-registered threshold $\kap \ge 0.25$.
Hypothesis H1 pre-registers that this boundary lies at the Q4\_K\_M$\to$Q3\_K\_M
transition for at least 2 of 3 model families.

\subsection{Why Existing Defenses Are Insufficient}

No existing prompt-injection defense system was designed for this failure mode in
an edge-hardware envelope.
Llama Guard~\citep{inan2023llama_guard}, PromptGuard~2, NeMo Guardrails,
Rebuff, SecAlign, Constitutional Classifiers~\citep{sharma2025cc},
LlamaFirewall~\citep{llamafirewall2025}, CaMeL, AEGIS~\citep{ghosh2024aegis}, and
SmoothLLM are all designed for server-side or full-precision deployment, are not
quantization-aware, and provide no formal portability guarantee across schemes.
\CG{} is, to our knowledge, the only system with all of: quantization-aware
calibration, FPR-decoupling theorem, online bandit adaptation, four-tier MVDS,
weight-hash attestation, and a black-box fallback path.

\subsection{Novel Adversary: A7 Quantization-Cliff Exploiter}

We introduce \textbf{A7}: an adversary who crafts natural-language prompts whose
refusal margin is just above $\tauq$ under FP16 but falls below it under Q3\_K\_M
/ NF4.
Because the FP16 model refuses, such prompts pass content moderation at upload
time; they succeed silently after quantization for edge deployment.
This adversary is enabled by the Egashira, Hong, and Mind the Gap results and has
no published defense.

% ─────────────────────────────────────────────────────────────────────────────
\section{System Architecture: Eight Components}
\label{sec:architecture}
% ─────────────────────────────────────────────────────────────────────────────

\CG{} comprises eight named components operating in a coordinated pipeline.

\subsection{VESTIBULE — Stateless Write-Side Input Gates}

VESTIBULE is a stack of fast, stateless gates over the raw prompt and any
retrieved/tool content:

\begin{description}[leftmargin=2em, itemsep=2pt]
  \item[VESTIBULE-LZ] LZ4/zlib compression ratio $\rho_\text{LZ}(x) = |z(x)|/|x|$.
        GCG adversarial suffixes are less compressible than English (near-random tokens);
        ArtPrompt~\citep{artprompt2024} is more compressible.
        Two-sided gate; ${\approx}20$\,\textmu s/KB, effectively free.
  \item[VESTIBULE-PS] Provenance-aware spotlighting~\citep{hines2024spotlighting}.
        Datamarks untrusted spans with a session-secret marker; asks the model to track
        provenance. Cuts injection ASR from $> 50\%$ to $< 2\%$ on capable models.
        On 2\,GB tiers: hard-delimited prompts only (datamarking omitted).
  \item[VESTIBULE-ID] Instruction-density gate: verb--entity collision counting to detect
        instruction hijacking without perplexity.
  \item[VESTIBULE-EA] Encoding-anomaly gate: non-printable character density, base-N
        detection, ASCII art density (targets A6).
\end{description}

VESTIBULE is stateless across requests; its configuration is identical across all
hardware tiers.

\begin{scaffoldbox}
VESTIBULE is implemented as Phase A scaffolding. All four sub-gates produce verdicts
on synthetic data. 939 tests pass. No real-traffic calibration has been run.
\end{scaffoldbox}

\subsection{PROBE — Quantization-Aware White-Box Observer}

PROBE attaches to the protected model's forward pass via registration hooks
(\texttt{register\_forward\_hook} in transformers; \texttt{cb\_eval} callback in
llama.cpp).
It does not introduce a separate classifier.

\subsubsection*{PROBE-RM (Refusal Margin)}

Pre-compute the refusal direction $\rhat_{M,q}$ via Arditi difference-in-means~\citep{arditi2024refusal}
at layer $\ell$ on a calibration set $\mathcal{C}$.
At inference, compute the \textbf{refusal margin}:

\begin{equation}
  \rho(x) = \frac{\langle h_\ell(x),\, \rhat \rangle}{\|\rhat\|}
  \label{eq:margin}
\end{equation}

at the post-instruction token $t_\text{post-inst}$, before generation.
Gate fires when $\rho(x) < \tauq$, where $\tauq$ is a per-$(M, q)$ quantile of the
benign calibration distribution at FPR $= \alpha = 0.05$.

\subsubsection*{PROBE-MT (Margin Trajectory)}

Computes $\rho_\ell$ at a small set of layers (e.g., 25\%/50\%/75\% depth).
Defines $\dot\rho = \rho_\text{late} - \rho_\text{mid}$.
Adversarial suffixes that redirect the model often show characteristic non-monotone
margin trajectories absent in benign prompts.

\subsubsection*{PROBE-HD (Harmfulness Direction)}

Following~\citet{zhao2025refusal}, projects the residual stream at the
\emph{user-instruction} token $t_\text{inst}$ onto the harmfulness direction $\hhat$.
Together with PROBE-RM this forms a 2-D safety subspace probe: $\rho_r$ catches
response-stage refusal, $\rho_h$ catches request-stage harmfulness recognition.

\subsubsection*{Quantization-Aware Calibration}

The calibration table $(\rhat_{M,q}, \tauq, \mu_0^{(q)})$ is recomputed once per
$(M, q)$ pair on a fixed calibration set $\mathcal{C}$.
This is the operational content of the FPR-decoupling theorem
(\S\ref{sec:decoupling}): only the calibration table is quantization-specific; the
gate code is portable.

\begin{testedbox}
PROBE-RM is tested with real model inference.
Run: \texttt{Llama-3.2-3B-Instruct}, FP16 and NF4, layer~14, Colab T4.
Calibrated thresholds: $\tauq_\text{FP16} = 0.09742$, $\tauq_\text{NF4} = 0.09827$
($\Delta\tauq = +0.00085$). 400 benign + 200 harmful prompts. Fold A complete.
\end{testedbox}

\begin{scaffoldbox}
PROBE-MT and PROBE-HD are implemented as Phase A scaffolding. They accept synthetic
hidden-state arrays and produce correct output shapes. Not tested with real inference.
\end{scaffoldbox}

\subsection{B-PROBE — Black-Box Fallback Observer}

For closed-weight API endpoints and NPU-frozen graphs without residual-stream access.

\subsubsection*{B-PROBE-LOGIT}

Let $\ell(q) = (\log p(v_1 \mid q), \ldots, \log p(v_k \mid q))$ be the top-$k$
log-probability vector at the first response token.
Train a logistic head $\varphi_\theta : \mathbb{R}^k \to \mathbb{R}$ on the
calibration set $\mathcal{C}$ with labels (refused/complied).
B-PROBE-LOGIT margin: $m_\text{BL}(q) = \varphi_\theta(\ell(q))$.

The first-token logit gap between \{``I cannot'', ``I'm sorry''\} and
\{``Sure'', ``Here''\} is a noisy linear functional of $\rhat^\top z$ via the
unembedding $W_U$.
FPR portability holds (Corollary~\ref{cor:blackbox}); TPR is strictly weaker and
pre-registered as Hypothesis H3.

\subsubsection*{B-PROBE-CONSISTENCY}

Generate $N \in \{2, 3\}$ paraphrases of the prompt using a \emph{local}
paraphraser (1B distilled, different family from test model).
Query the closed endpoint.
Compute pairwise Jensen--Shannon divergence of first-token distributions.
High divergence indicates the prompt is on the cliff; small perturbations flip behavior.

\begin{scaffoldbox}
Both B-PROBE sub-primitives are Phase A scaffolding. Top-k logprob integration with
a real API or inference engine has not been tested.
\end{scaffoldbox}

\subsection{TRIPWIRE — Streaming Change-Point and Entropy Monitor}

TRIPWIRE operates \emph{during decoding} without buffering outputs.

\subsubsection*{TRIPWIRE-H (Per-Token Entropy CUSUM)}

Per-token entropy $H_t = -\sum_v p_t(v) \log p_t(v)$ over the model's top-$k$
softmax.
One-sided CUSUM on entropy residuals:

\begin{equation}
  S_t = \max\!\left(0,\; S_{t-1} + (H_t - \mu_0 - \delta/2)\right),
  \quad \text{alarm if } S_t > h.
  \label{eq:cusum}
\end{equation}

Threshold $h$ tuned for $\mathrm{ARL}_0 \approx 1000$ benign tokens
(roughly one false alarm per 50--100 sessions at typical conversation lengths).
$\mu_0$ is the per-$(M, q)$ benign mean estimated on the calibration fold.

An injection pushing the model toward a forbidden completion typically shows
an entropy depression (high-confidence forbidden token) followed by an entropy spike
(model ``starts and stalls'').
GCG-perturbed contexts often show abnormal early entropy; many-shot injections show
monotonic NLL decrease detectable by the CUSUM.

\subsubsection*{TRIPWIRE-R (KenLM Reference-Ratio Gate)}

Train a KenLM 5-gram model offline on a benign instruction corpus.
At decode time: $\Delta_t = \log P_M(y_t \mid \mathrm{ctx}) - \log P_\text{ref}(y_t \mid \mathrm{ctx})$.
This approximates the Neyman--Pearson log-likelihood ratio.
Combined with TRIPWIRE-H via CUSUM over $\alpha H_t + \beta \Delta_t$ where
$(\alpha, \beta)$ is fit on calibration traces.
Because the reference $n$-gram is fixed, quantization shifts only the $\log P_M$
term; the ratio absorbs scheme-specific variation better than bare perplexity.

\begin{scaffoldbox}
Both TRIPWIRE sub-primitives are Phase A scaffolding. Streaming integration with
real model decoding has not been tested. KenLM is not yet trained on any corpus.
\end{scaffoldbox}

\subsection{CONDUCTOR — LinUCB Contextual-Bandit Gate Orchestrator}

CONDUCTOR selects a gate-weight vector per request.
It is the only learnable component; all others are calibration-only.

\subsubsection*{Formulation}

\textbf{Arms.} A finite menu $\mathcal{A}$ of preset gate-weight vectors, $|\mathcal{A}| = 8$--16.

\textbf{Context.} $x_t \in \mathbb{R}^d$ with $d \approx 32$: prompt length, character entropy,
language ID, has-tool-context flag, hardware tier, recent-incident EWMA, TRIPWIRE-H
baseline drift estimate.
\textbf{No payload content is stored.}

\textbf{Policy.} LinUCB~\citep{li2010contextual}:

\begin{equation}
  a_t = \arg\max_a \left( x_t^\top \hat\theta_a +
        \alpha \sqrt{x_t^\top A_a^{-1} x_t} \right),
  \label{eq:linucb}
\end{equation}

with $A_a = \lambda I + \sum_{s : a_s = a} x_s x_s^\top$,
$b_a = \sum r_s x_s$, $\hat\theta_a = A_a^{-1} b_a$.
Regret: $\tilde O(\sqrt{dT})$ in the stochastic regime.

\textbf{Adversarial fallback.} EXP3.S~\citep{auer2002exp3} when incident-rate EWMA
crosses a threshold, giving $\tilde O(\sqrt{KTS})$ regret for $S$ distribution
switches.

\textbf{Pre-registered reward function.}
\begin{equation}
  r_t = \mathbb{1}[\text{served} \land \text{clean}]
       - \mathbb{1}[\text{served} \land \text{injected}]
       - 0.2 \cdot \mathbb{1}[\text{blocked} \land \text{benign}].
  \label{eq:reward}
\end{equation}

Updates on: LOOKOUT verdict, canary-token leakage, tool-call audit failure,
user complaint flag.

\textbf{Cold-start.} Fresh deployments use the static maximum-coverage policy until
$O(d)$ requests per arm accrue.

\textbf{Safe rollback (A8 defense).}
If the cumulative ABR over $W$ rounds exceeds a pre-registered $A_\text{max}$ for
two consecutive windows, revert to the static maximum-coverage policy for $K$ rounds,
then resume exploration.

\begin{scaffoldbox}
CONDUCTOR is Phase A scaffolding. LinUCB and EXP3.S are implemented and accept
synthetic reward signals. No live deployment or bandit training has been run.
Fold D (bandit drift) is planned.
\end{scaffoldbox}

\subsection{LOOKOUT — Output-Side Judge and Canary Loop}

\subsubsection*{LOOKOUT-CT (Canary Token Loop)}

Per-session CSPRNG-generated canary token sequence inserted into the system prompt.
LOOKOUT scans output streams (via a Count-Min sketch keyed by the session secret)
for any echo.
Canary is per-session, never logged, discarded at session end.
Bloom filter: $m = 256$ bits, $k = 3$ hashes; FPR $\approx 0.001$ at typical
canary lengths.
Total session state: $\approx 32$ bytes.

\subsubsection*{LOOKOUT-JG (KL-Drift / Mutation Judge)}

Generate $N \in \{2, 3\}$ character-level mutations of the input.
Take only the first-token distributions of the model on each.
Compute pairwise JS divergence.
High divergence on the \emph{first token} indicates an unstable input---
characteristic of GCG/AutoDAN-perturbed prompts.
Cost: $N$ extra prefill calls of one token each (not full responses).
On 8\,GB tier: $N=3$; Pi 5: $N=2$; 2\,GB tier: omitted.

\subsubsection*{LOOKOUT Classifier}

On 8\,GB tier: Llama-Guard-3-1B-INT4~\citep{inan2023llama_guard} (440\,MB) or
DeBERTa-86M.
On Pi 5: DeBERTa-22M or PromptGuard-2-22M~\citep{llamafirewall2025}.
On 2\,GB tier: rules only.
Provides the reward signal $r_t$ to CONDUCTOR.

\begin{scaffoldbox}
All LOOKOUT sub-components are Phase A scaffolding. Canary logic, JG, and classifier
wrappers accept synthetic data. No real-inference integration tested.
\end{scaffoldbox}

\subsection{LADDER — Four-Tier MVDS Hardware Router}

LADDER is configuration, not learning. Selects which gates run at what depth based
on detected hardware and observability.

\begin{table}[H]
  \centering
  \small
  \caption{Four hardware tiers and their gate configurations.}
  \label{tab:tiers}
  \begin{tabular}{@{}lp{2.8cm}p{2.2cm}p{4.2cm}l@{}}
    \toprule
    Tier & Hardware & Schemes & Active Gates & Scope \\
    \midrule
    \textbf{A} & RTX 5060 8\,GB & FP16, NF4, AWQ-INT4, Q4--Q6 & All 12 gates (incl. LOOKOUT-JG) & Full stack, LinUCB $|\mathcal{A}|=16$ \\[3pt]
    \textbf{B} & Raspberry Pi 5 8\,GB CPU & GGUF Q4\_K\_M, Q3\_K\_M & 11 gates (LOOKOUT-JG $N=2$) & LinUCB $|\mathcal{A}|=8$; B-PROBE-CONSISTENCY subs for JG \\[3pt]
    \textbf{C} & 2\,GB embedded (RK3588, Jetson Orin, Pi 4) & Q3\_K\_M, Q2\_K, RKNN W8A8 & VESTIBULE-LZ, PROBE-RM (1 layer), Page--Hinkley, LOOKOUT-CT & \textbf{No bandit}; static weights; \textbf{NOT defended against A7} \\[3pt]
    \textbf{C+} & 2\,GB + PromptGuard-2-22M-INT4 & Q3\_K\_M, RKNN W8A8 & VESTIBULE-LZ/PS, B-PROBE-LOGIT, PROBE-RM, LOOKOUT-CT, PG2 & Static weights; H5 tests this tier \\
    \bottomrule
  \end{tabular}
\end{table}

\noindent\textbf{Plain-language statement on Tier C:}
Tier C runs Q3\_K\_M models in 2\,GB---the very regime where Egashira and Hong
et al.\ cliffs are most likely.
The minimal gate stack raises attacker cost only marginally.
\textbf{Tier C is NOT FOR OPEN-DOMAIN ADVERSARIAL USE.}
Suitable only for single-task Jatmo-style~\citep{jatmo2023} assistants with a fixed
input grammar and no agentic tool use.

\begin{scaffoldbox}
LADDER is implemented as a configuration-based tier dispatcher. All tier routing
logic is scaffolded and tested with synthetic data.
\end{scaffoldbox}

\subsection{ATTEST — Boot-Time Weight-Hash Attestation}

Boot-time SHA-256 of the GGUF/safetensors blob compared against a vendor-signed
manifest.
Defends against A2 (Egashira-style poisoned weights~\citep{egashira2024exploiting}).

Llama, Qwen, Mistral, Gemma publish per-file SHA-256 on Hugging Face.
For unsigned community quantizations, ATTEST falls back to first-use trust with
a sticky local hash.

\begin{equation}
  \text{ATTEST}(M) =
  \begin{cases}
    \text{ALLOW}    & \text{if } \mathrm{SHA256}(\text{blob}) = \text{manifest hash} \\
    \text{BLOCK}    & \text{if mismatch}
  \end{cases}
\end{equation}

\begin{scaffoldbox}
ATTEST-WH is Phase A scaffolding. The SHA-256 hashing and manifest comparison logic
are implemented. No live boot-time integration with an inference engine has been tested.
\end{scaffoldbox}

% ─────────────────────────────────────────────────────────────────────────────
\section{Theoretical Foundations}
\label{sec:theory}
% ─────────────────────────────────────────────────────────────────────────────

\subsection{Geometric Assumption}

\begin{definition}[Safety Subspace, G1]
Let $z_t(q)$ denote the residual-stream activation at token $t$ for prompt $q$.
There exists a subspace $S = \mathrm{span}\{\rhat, \hhat\} \subset \mathbb{R}^d$
such that for any distribution $P$ over prompts,
\[
\|\Pi_S z_{t_\text{inst}}(q) - \Pi_S z_{t_\text{inst}}(q')\|
\;\ge\; \xi \cdot \|z_{t_\text{inst}}(q) - z_{t_\text{inst}}(q')\|
\]
in expectation over $(q, q') \in \mathcal{H} \times \mathcal{B}$, for some $\xi > 0$.
\end{definition}

G1 is empirically supported by~\citet{arditi2024refusal} and~\citet{zhao2025refusal}
for instruction-tuned models at 7B--72B parameters.

\textbf{Sensitivity bound.}
If quantization $\tilde q$ acts as a Lipschitz operator on the residual stream with
constant $L_{\tilde q}$, then
\begin{equation}
  \|\pi(q) - \pi(q')\| \;\le\; L_{\tilde q} \cdot (\|\delta_z\|_2 + 2\varepsilon_{\tilde q}),
  \label{eq:sensitivity}
\end{equation}
where $\varepsilon_{\tilde q}$ is the per-quantization residual reconstruction error
and $\delta_z$ is the induced activation perturbation.
The \textbf{cliff is the regime where $\varepsilon_{\tilde q}$ inflates enough that
$\pi(q)$ flips sign on harmful prompts}.

\subsection{FPR-Decoupling Theorem}
\label{sec:decoupling}

\begin{theorem}[FPR Decoupling, White-Box]
\label{thm:decoupling}
Let $g(x) = \mathbb{1}[\phi(x; M_q) > \tau_q]$ be a write-side gate operating on
(a) input strings and/or (b) summary statistics of model outputs (top-$k$
log-probabilities, refusal-margin scalars).
Let $\tau_q$ be the empirical $(1-\alpha)$-quantile of $\phi(\cdot; M_q)$ on a
calibration set $\mathcal{C} \sim \mathcal{B}$.
Then
\[
|\mathrm{FPR}_{g, q} - \mathrm{FPR}_{g, q^*}| \;\le\; \mathrm{KS}(F_q, F_{q^*}) \;\le\; d_{q, q'}
\]
as $|\mathcal{C}| \to \infty$, where $d_{q, q'}$ is the Kolmogorov--Smirnov distance
between benign-margin distributions.
\end{theorem}

\begin{proof}[Proof sketch]
By Glivenko--Cantelli, the empirical quantile of $\phi(\cdot; M_q)$ on $\mathcal{C}$
converges uniformly to the population quantile at rate $O(|\mathcal{C}|^{-1/2})$,
regardless of the distribution of $\phi$.
Thus the realized FPR of each gate is $\alpha + o(1)$ at any fixed $\alpha$.
Composition by AND/OR/weighted-sum is a finite function of finitely many
$\alpha$-controlled gates; the result follows.
The KS bound follows from the definition of CDF inverse and KS distance.
\end{proof}

\begin{corollary}[FPR Decoupling, Black-Box]
\label{cor:blackbox}
The same theorem applies to B-PROBE-LOGIT's logistic head $\varphi_\theta : \mathbb{R}^k \to \mathbb{R}$
trained on (prompt, refused/complied) pairs.
TPR is \emph{not} decoupled in either case.
\end{corollary}

\noindent\textbf{Honest scope.}
Theorem~\ref{thm:decoupling} protects FPR only.
It says nothing about TPR.
The cliff is \emph{defined} as the regime where TPR collapses while FPR is maintained.

\noindent\textbf{Sensitivity corollary.}
With $|\mathcal{C}| < 100$, the KS estimation error dominates, giving an $O(1/\sqrt{|\mathcal{C}|})$
inflation of the bound.
Deployments with $|\mathcal{C}| = 2000$ from Anthropic-HH suffice; deployments with
$|\mathcal{C}| < 100$ must accept a wider FPR portability band.

\subsection{Information-Theoretic Framing}

TRIPWIRE approximates the Neyman--Pearson log-likelihood ratio
$\Lambda = \log P_1/P_0$, which CUSUM (Page 1954) accumulates with optimal expected
detection delay (Lorden 1971, Moustakides 1986).
Pure perplexity gates approximate $\log P_0$ alone; TRIPWIRE approximates the ratio
by combining a fixed KenLM reference with the model's own per-step entropy.

\subsection{RL-Theoretic Framing}

CONDUCTOR is a contextual adversarial bandit.
LinUCB gives $\tilde O(\sqrt{dT})$ regret in the stochastic regime.
Under non-stationary attack distributions, EXP3.S~\citep{auer2002exp3} gives
$\tilde O(\sqrt{KTS})$ for $S$ switches---appropriate during coordinated attack
campaigns.

% ─────────────────────────────────────────────────────────────────────────────
\section{Threat Model (A1--A9)}
\label{sec:threats}
% ─────────────────────────────────────────────────────────────────────────────

We adopt a Greshake-style~\citep{greshake2023not} hierarchy expanded for the
edge-quantized setting.
No persistent storage of user payloads; only ephemeral session windows, RL policy
weights, sketch counters, and signal buffers may be retained.

\begin{table}[H]
  \centering
  \small
  \caption{Threat model: nine adversaries, capabilities, and primary \CG{} defenses.
           A7 is the novel adversary introduced by this work.}
  \label{tab:threats}
  \begin{tabularx}{\textwidth}{@{}lp{4.0cm}p{3.0cm}X@{}}
    \toprule
    ID & Adversary & Capability & Primary Defense \\
    \midrule
    A1 & Direct injector & NL hijacks, DAN, role-play, persuasion & VESTIBULE, LOOKOUT \\
    A2 & Indirect injector / poisoned-weight attacker~\citep{egashira2024exploiting} & Embeds instructions in RAG; FP16-benign / Q3\_K\_M-malicious weights & VESTIBULE-PS, ATTEST + LADDER \\
    A3 & Optimizer~\citep{zou2023advbench} & GCG, AutoDAN, AmpleGCG adversarial suffixes & PROBE, VESTIBULE-LZ \\
    A4 & Iterator & PAIR, TAP, Crescendo~\citep{crescendo2024} black-box query loops & TRIPWIRE, LOOKOUT \\
    A5 & Scaler & Best-of-N~\citep{bon2024}, many-shot, sampling variance & TRIPWIRE-H, LOOKOUT-CT \\
    A6 & Encoder & ArtPrompt~\citep{artprompt2024}, cipher, low-resource language & TRIPWIRE-R \\
    \textbf{A7} & \textbf{Quantization-cliff exploiter} & Natural-language prompts that refuse at FP16 but comply at Q3\_K\_M / NF4 & LADDER + ATTEST + cross-tier rollback \\
    A8 & Defender-aware adversary & Knows policy weights and calibration tables (Kerckhoffs assumption) & CONDUCTOR safe-rollback \\
    A9 & Closed-weight black-box endpoint adversary & Targets API endpoint; only text and top-$k$ logprobs observable & B-PROBE + LOOKOUT-JG \\
    \bottomrule
  \end{tabularx}
\end{table}

% ─────────────────────────────────────────────────────────────────────────────
\section{Cliff Metrics: Formal Definitions}
\label{sec:metrics}
% ─────────────────────────────────────────────────────────────────────────────

\subsection{Geometric Cliff $\cliff$}

Let $M_q$ be the model quantized to scheme $q$ and $\mathcal{H}$ a held-out harmful
set.
Define the \textbf{separation} at scheme $q$:
\[
\Delta_q = \mathrm{med}(\mathcal{R}_q^\mathcal{B}) - \mathrm{med}(\mathcal{R}_q^\mathcal{H}),
\]
where $\mathcal{R}_q^\mathcal{X}$ is the empirical distribution of $\rho_q(x)$
for $x \in \mathcal{X}$.
The \textbf{safety-cliff metric}:
\begin{equation}
  \cliff(q; q^*) = \frac{\Delta_{q^*} - \Delta_q}{\Delta_{q^*}}.
  \label{eq:cliff_formal}
\end{equation}
A \textbf{safety cliff} between $q_1$ and $q_2$ (more aggressive) exists when
$\cliff(q_2; q^*) - \cliff(q_1; q^*) \ge \kap$, for $\kap = 0.25$.

\subsection{Wasserstein Cliff $\wcliff$}

\begin{equation}
  \wcliff(M, q; q^*) = W_2(P_q, P_{q^*}),
  \label{eq:wcliff}
\end{equation}
where $P_q$ is the empirical distribution of $\rho_q(x)$ over $\mathcal{H}$.
Captures lower-tail movement (operationally dangerous tail) that the median misses.

\subsection{Behavioral Cliff $\bcliff$}

\begin{equation}
  \bcliff(M, q; q^*) = |C_{M,q}(\mathcal{H}) - C_{M,q^*}(\mathcal{H})|,
  \label{eq:bcliff}
\end{equation}
where $C_{M,q}(\mathcal{H})$ is the compliance rate under the judge stack
$J = \text{StrongREJECT}$~\citep{souly2024strongreject} $\land$
Llama-Guard-3-8B (both must label the response as harmful-and-substantively-helpful).
This metric is \textbf{non-circular}: it does not use PROBE-RM margins.

\subsection{H1 Criterion}

H1 is accepted for model family $\mathcal{F}$ if and only if both
$\cliff(q^*) \ge \kap$ AND $\bcliff(q^*) \ge \kap$ at the same scheme boundary,
for $\ge 2$ of 3 model families.
\textbf{Disagreement between geometric and behavioral metrics falsifies the
PROBE-RM design center.}

% ─────────────────────────────────────────────────────────────────────────────
\section{Five Pre-Registered Hypotheses}
\label{sec:hypotheses}
% ─────────────────────────────────────────────────────────────────────────────

\begin{table}[H]
  \centering
  \small
  \caption{Five pre-registered hypotheses, acceptance criteria, and empirical status.}
  \label{tab:hypotheses}
  \begin{tabularx}{\textwidth}{@{}lp{3.8cm}p{3.2cm}lX@{}}
    \toprule
    & Hypothesis & Claim & Threshold & \textbf{Status} \\
    \midrule
    H1 & Cliff existence & $\cliff$ and $\bcliff$ jump $\ge\kap$ at same boundary, $\ge 2/3$ families & $\kap = 0.25$ & \textcolor{orange!70!black}{\textbf{Preliminary.}} 2 schemes, 1 family. H1 not accepted. Q3\_K\_M untested. \\[4pt]
    H2 & FPR decoupling (white-box) & PROBE-RM FPR varies $< \varepsilon$ across 5 schemes after calibration & $\varepsilon = 0.02$ & \textcolor{orange!70!black}{\textbf{Preliminary signal}} ($\Delta\tauq = 0.00085$ FP16$\to$NF4). Needs 5 schemes. \\[4pt]
    H3 & FPR decoupling (black-box) & B-PROBE-LOGIT FPR varies $< \varepsilon$; TPR strictly $<$ PROBE-RM & $\varepsilon = 0.02$ & \textcolor{red!60!black}{\textbf{Not tested.}} B-PROBE not integrated. \\[4pt]
    H4 & Composition gain & Full stack ABR $<$ any single primitive at matched FPR & Wilcoxon $p < 0.01$ & \textcolor{red!60!black}{\textbf{Not tested.}} Fold C not run. \\[4pt]
    H5 & Tier-C structural weakness & Tier C: no significant ABR gain vs.\ no-defense; Tier C+ does & Wilcoxon $p$(C) $\ge 0.05$; $p$(C+) $< 0.05$ & \textcolor{red!60!black}{\textbf{Not tested.}} Tier C deployment not evaluated. \\
    \bottomrule
  \end{tabularx}
\end{table}

% ─────────────────────────────────────────────────────────────────────────────
\section{Five-Fold Evaluation Protocol}
\label{sec:folds}
% ─────────────────────────────────────────────────────────────────────────────

\begin{table}[H]
  \centering
  \small
  \caption{Five evaluation folds, datasets, outputs, and current status.}
  \label{tab:folds}
  \begin{tabularx}{\textwidth}{@{}llp{4.0cm}p{3.0cm}X@{}}
    \toprule
    Fold & Purpose & Constructed from & Computes & \textbf{Status} \\
    \midrule
    \textbf{A} & Calibration & Anthropic-HH, OASST1 (2K benign + 2K refused) & $\rhat_{M,q}$, $\tauq$, $\mu_0$, KenLM 5-gram & \textcolor{green!50!black}{\textbf{TESTED}} (400 benign + 200 harmful; Llama-3.2-3B; FP16 + NF4; layer 14) \\[4pt]
    \textbf{B} & Cliff measurement & AdvBench~\citep{zou2023advbench} + HarmBench~\citep{mazeika2024harmbench} (held out from A) & $\cliff$, $\wcliff$, $\bcliff$ & \textcolor{green!50!black}{\textbf{TESTED}} ($\bcliff$ uses PROBE-RM proxy, not full StrongREJECT+LG3-8B) \\[4pt]
    \textbf{C} & Defense composition & JBB~\citep{chao2024jailbreakbench}, AgentDojo, InjecAgent, GCG, AutoDAN, PAIR, TAP, Crescendo, ArtPrompt, BoN, TensorTrust & per-primitive ABR, FPR; composition curves & \textcolor{red!60!black}{\textbf{NOT RUN}} \\[4pt]
    \textbf{D} & Bandit / online drift & Streaming mixture from Fold C with synthetic distribution shifts & CONDUCTOR regret, rollback frequency & \textcolor{red!60!black}{\textbf{NOT RUN}} \\[4pt]
    \textbf{E} & BCN-2 dataset construction & Paraphrases generated by \textbf{Mistral-7B-base} (non-RLHF, non-test-family) of AdvBench, filtered by FP16 refusal on test family & Held-out BCN-2 corpus for A7 Fold C evaluation & \textcolor{red!60!black}{\textbf{NOT RUN}} \\
    \bottomrule
  \end{tabularx}
\end{table}

\subsection{BCN-2 Anti-Circularity Protocol}

The central anti-circularity move: the dataset's \emph{construction filter} uses the
FP16 model's behavior; the \emph{paraphrasing engine} is from a different family
than the test family; the cliff being tested is a property of the test family at
low bit-width, independent of the construction filter.

Three pre-registered alternatives:
\begin{enumerate}[leftmargin=2em]
  \item \textbf{BCN-2 (primary):} Mistral-7B-base paraphraser; FP16-refusal filter on test family.
  \item \textbf{Mechanistic-interpretability-derived:} Use $\hhat$ (Zhao et al.) at FP16 to
        identify prompts with high harmfulness projection but borderline refusal projection.
  \item \textbf{Distribution-shift-based:} Cliff as $W_2(P^\text{FP16}, P^{\text{Q3\_K\_M}})$
        on the first-token output distribution over a benign corpus.
\end{enumerate}

\subsection{Statistical Power}

For H1 with $\kap = 0.25$, $\alpha = 0.05$, power $0.8$, two-sided two-sample
$t$-test on median margins, with $\sigma \approx 0.18$ (from Hong et al.):

\[
n_\text{min} = \frac{2 \cdot (z_{1-\alpha/2} + z_{1-\beta})^2 \cdot \sigma^2}{\kap^2}
\approx \frac{2 \cdot (1.96 + 0.84)^2 \cdot 0.0324}{0.0625} \approx 8.1.
\]

We pre-register $n = 200$ per $(M, q)$ cell ($\gg n_\text{min}$) to absorb prompt-difficulty
stratification and multilingual subgroups.
Total Fold B sample: 3 families $\times$ 6 quantizations $\times$ 200 = 3\,600 prompts.

% ─────────────────────────────────────────────────────────────────────────────
\section{Empirical Results: Folds A and B}
\label{sec:results}
% ─────────────────────────────────────────────────────────────────────────────

\subsection{Setup}

\textbf{Model:} \texttt{meta-llama/Llama-3.2-3B-Instruct}, layer $\ell = 14$ (of 28).\\
\textbf{Hardware:} Google Colab T4 (NVIDIA Tesla T4, 16\,GB VRAM, CUDA 12.8).\\
\textbf{Software:} \texttt{transformers} 4.48, \texttt{bitsandbytes} 0.45.49.2,
\texttt{torch} 2.5, \texttt{numpy} 1.26 ($< 2$ required by bitsandbytes).\\
\textbf{Schemes tested:} FP16 ($\approx 6.2$\,GB VRAM), NF4 ($\approx 2.1$\,GB VRAM).\\
\textbf{Schemes \emph{not} tested:} INT8, AWQ-INT4, GGUF Q4\_K\_M, Q3\_K\_M, Q2\_K
(the predicted cliff zone).

\subsection{Fold A: Calibration Results}

\begin{testedbox}
400 benign prompts (Anthropic-HH-RLHF helpful turns + OASST1) and 200 harmful prompts
(Anthropic-HH-RLHF rejected turns).
FPR target $\alpha = 0.05$.
\end{testedbox}

\begin{table}[H]
  \centering
  \caption{Fold A calibration: PROBE-RM refusal thresholds per scheme.}
  \label{tab:fold_a}
  \begin{tabular}{lccc}
    \toprule
    Scheme & $\tauq$ & $\Delta\tauq$ vs.\ FP16 & Layer \\
    \midrule
    FP16 & $0.09742$ & --- & 14 \\
    NF4  & $0.09827$ & $+0.00085$ ($+0.87\%$) & 14 \\
    \bottomrule
  \end{tabular}
\end{table}

The NF4 threshold exceeds FP16 by 0.87\%.
The small $\Delta\tauq$ is a preliminary positive signal for H2 (FPR decoupling):
calibration absorbs most of the quantization-induced margin shift in the benign
distribution.
With only two schemes, this is insufficient to accept H2.

\subsection{Fold B: Cliff Measurement Results}

\begin{testedbox}
Fold B corpus: AdvBench~\citep{zou2023advbench} (520 harmful behaviors) +
JailbreakBench~\citep{chao2024jailbreakbench} (100 harmful goals, JBB-Behaviors).
Fold isolation verified via SHA-256 disjointness before execution.
$\bcliff$ uses PROBE-RM margin threshold crossing as behavioral proxy (not full
StrongREJECT + Llama-Guard-3-8B judge --- see limitation below).
\end{testedbox}

\begin{table}[H]
  \centering
  \caption{Fold B cliff measurement results: Llama-3.2-3B-Instruct, Layer 14.
           $\dagger$~$\bcliff$ uses PROBE-RM proxy (partially circular); the
           full StrongREJECT + Llama-Guard-3-8B judge is Phase C work.}
  \label{tab:fold_b}
  \begin{tabular}{lcccc}
    \toprule
    Scheme & $\cliff$ & $\wcliff$ & $\bcliff{}^\dagger$ & H1 crossed ($\ge\kap=0.25$)? \\
    \midrule
    FP16 & 0.000 & 0.000 & 0.000 & --- (baseline) \\
    NF4  & \textbf{0.167} & 0.014 & 0.000 & \textbf{No} \\
    \midrule
    \multicolumn{5}{l}{\texttt{cliff\_boundary = null} \quad \texttt{h1\_accepted = false}} \\
    \bottomrule
  \end{tabular}
\end{table}

\paragraph{Geometric metric ($\cliff = 0.167$).}
The NF4 refusal direction has rotated from the FP16 baseline.
At $0.67\kap$, this is a substantive sub-threshold signal.
Consistent with the expectation that NF4's normal-quantile design
preserves activation statistics better than integer quantization.

\paragraph{Wasserstein metric ($\wcliff = 0.014$).}
Near zero: the margin \emph{distribution} over benign prompts is largely unchanged
even as the direction has rotated.

\paragraph{Behavioral proxy ($\bcliff = 0.000$).}
No additional harmful prompts crossed the calibrated threshold under NF4.
\textbf{Critical limitation:} the proxy uses the same PROBE-RM margin threshold that
was calibrated in Fold A---partially circular.
A threshold calibrated on the benign distribution will not register sub-threshold
shifts as behavioral changes.
The full StrongREJECT + Llama-Guard-3-8B judge is required for a non-circular verdict.

\paragraph{Geometric--behavioral discordance.}
$\cliff = 0.167$ versus $\bcliff = 0.000$ (proxy) is a discordant pair.
Under the full judge, $\bcliff$ may remain near zero (NF4 shift too small for
actual behavioral change) or may register a non-zero value.
This is the primary motivation for Fold C.

% ─────────────────────────────────────────────────────────────────────────────
\section{Implementation Status Summary}
\label{sec:status}
% ─────────────────────────────────────────────────────────────────────────────

\begin{table}[H]
  \centering
  \small
  \caption{Full implementation and testing status as of May 2026.}
  \label{tab:status}
  \begin{tabularx}{\textwidth}{@{}lXll@{}}
    \toprule
    Component / Feature & Details & Implemented & Tested \\
    \midrule
    VESTIBULE-LZ & LZ compression gate & Phase A stub & Synthetic only \\
    VESTIBULE-PS & Spotlighting gate & Phase A stub & Synthetic only \\
    VESTIBULE-ID, -EA & Density + encoding gates & Phase A stub & Synthetic only \\
    \midrule
    PROBE-RM & Difference-in-means + margin & \YES{} Full & \YES{} Real (Fold A) \\
    PROBE-MT & Multi-layer trajectory & Phase A stub & Synthetic only \\
    PROBE-HD & Harmfulness direction & Phase A stub & Synthetic only \\
    \midrule
    B-PROBE-LOGIT & Top-$k$ logprob head & Phase A stub & Synthetic only \\
    B-PROBE-CONSISTENCY & $N$-paraphrase JSD & Phase A stub & Synthetic only \\
    \midrule
    TRIPWIRE-H & Entropy CUSUM & Phase A stub & Synthetic only \\
    TRIPWIRE-R & KenLM reference-ratio & Phase A stub & Synthetic only \\
    \midrule
    CONDUCTOR & LinUCB + EXP3.S & Phase A stub & Synthetic only \\
    \midrule
    LOOKOUT-CT & Canary token loop & Phase A stub & Synthetic only \\
    LOOKOUT-JG & Mutation JSD judge & Phase A stub & Synthetic only \\
    LOOKOUT classifier & LG3/DeBERTa wrapper & Phase A stub & Synthetic only \\
    \midrule
    LADDER & Tier routing & Phase A stub & Synthetic only \\
    \midrule
    ATTEST-WH & SHA-256 manifest check & Phase A stub & Synthetic only \\
    \midrule
    Fold A (calibration) & PROBE-RM on real data & \YES{} & \YES{} (FP16, NF4) \\
    Fold B (cliff metrics) & 3 metrics, 2 schemes & \YES{} & \YES{} ($\bcliff$ proxy) \\
    Fold C (defense composition) & ABR/FPR per gate & \NO{} & \NO{} \\
    Fold D (bandit drift) & CONDUCTOR live & \NO{} & \NO{} \\
    Fold E (BCN-2 dataset) & Non-circular construction & \NO{} & \NO{} \\
    \midrule
    GGUF Q4\_K\_M / Q3\_K\_M / Q2\_K & llama.cpp integration & \NO{} & \NO{} \\
    AWQ-INT4 & autoawq integration & \NO{} & \NO{} \\
    Multiple model families & Qwen-2.5, Mistral, Gemma & \NO{} & \NO{} \\
    Full $\bcliff$ judge & StrongREJECT + LG3-8B & \NO{} & \NO{} \\
    \bottomrule
  \end{tabularx}
\end{table}

\textbf{Test suite:} 939 tests passing, mypy strict on 53 source files, ruff clean.
All Phase A scaffolding tests operate on synthetic data and verify API contracts,
not real model behavior.

% ─────────────────────────────────────────────────────────────────────────────
\section{Quantization-Impact Analysis}
\label{sec:quant_impact}
% ─────────────────────────────────────────────────────────────────────────────

\begin{table}[H]
  \centering
  \small
  \caption{Per-scheme expected safety impact, based on published evidence.
           \emph{Values in the ``Measured'' column are from this work.}
           All other values are predictions derived from cited literature.}
  \label{tab:quant_impact}
  \begin{tabularx}{\textwidth}{@{}lXXXlX@{}}
    \toprule
    Scheme & Expected $\cliff$ & FPR shift & Recalibration & Source & \textbf{Measured} \\
    \midrule
    FP16 & 0.000 (reference) & reference & n/a & --- & \textbf{0.000} \\
    INT8 (LLM.int8()) & small & small & low & bitsandbytes docs & --- \\
    NF4 (bitsandbytes) & moderate & moderate & medium & \citet{hong2024quant} & \textbf{0.167} \\
    AWQ-INT4 & small--moderate & small--moderate & low--medium & Lin et al.\ (2024) & --- \\
    GGUF Q6\_K, Q5\_K\_M & small & small & low & \citet{mindthegap2025} & --- \\
    GGUF Q4\_K\_M & moderate & moderate & medium & \citet{egashira2024exploiting} & --- \\
    GGUF Q3\_K\_M & \textbf{large (cliff candidate)} & large & high & \citet{egashira2024exploiting}: ASR $\Delta = 88.7\%$ & \textbf{NOT TESTED} \\
    GGUF Q2\_K & catastrophic & unstable & gate may be unusable & Community reports & --- \\
    RKNN W8A8 & mixed & mixed & per-channel & Rockchip docs & --- \\
    \bottomrule
  \end{tabularx}
\end{table}

% ─────────────────────────────────────────────────────────────────────────────
\section{Comparison to Existing Defenses}
\label{sec:comparison}
% ─────────────────────────────────────────────────────────────────────────────

\begin{table}[H]
  \centering
  \small
  \caption{Comparison of \CG{} to eleven existing defense systems.
           WB = white-box (hidden state access); BB = black-box;
           QA = quantization-aware; Edge = edge-feasible;
           OA = online adaptive; FPR = formal FPR guarantee; SL = stateless.}
  \label{tab:comparison}
  \begin{tabularx}{\textwidth}{@{}lXccccccc@{}}
    \toprule
    System & WB/BB & QA & Edge & OA & FPR & SL & Source \\
    \midrule
    ContextForge & partial & \NO & partial & \NO & none & \YES & --- \\
    Constitutional Classifiers v1 & black & \NO & partial & \NO & none & server & \citet{sharma2025cc} \\
    Constitutional Classifiers++ & white (probe) & \NO & no (full model) & \NO & none & server & CC++ (2025) \\
    Llama Guard 3-1B-INT4 & black & indirect & \YES & \NO & none & \YES & \citet{inan2023llama_guard} \\
    PromptGuard 2-22M/86M & black & n/a & \YES & \NO & none & \YES & \citet{llamafirewall2025} \\
    LlamaFirewall & mixed & \NO & partial & partial & none & \YES & \citet{llamafirewall2025} \\
    NeMo Guardrails & depends & \NO & rules & none & \YES & --- & NVIDIA \\
    Rebuff & \YES & \NO & vector-DB & partial & none & partial & --- \\
    AEGIS & black & \NO & multi-LLM & \YES{} (Hedge) & no-regret & \YES & \citet{ghosh2024aegis} \\
    SmoothLLM & black & \NO & high cost & \NO & rand.\ smooth & n/a & SmoothLLM \\
    CaMeL & black & \NO & mid & \NO & info-flow & \YES & --- \\
    \midrule
    \textbf{\CG{}} & \textbf{both} & \textbf{\YES} & \textbf{\YES (A/B/C/C+)} & \textbf{\YES} & \textbf{decoupling} & \textbf{\YES} & \textit{this work} \\
    \bottomrule
  \end{tabularx}
\end{table}

\textbf{Distinct novelty.}
\CG{} is (1) the only system simultaneously quantization-aware and online-adaptive;
(2) the only system with a formal FPR-decoupling theorem and a black-box corollary;
(3) the only system with a non-circular cross-family cliff dataset protocol (BCN-2);
(4) the only system with explicit poisoned-weight attestation paired with runtime gates.

% ─────────────────────────────────────────────────────────────────────────────
\section{Per-Engine Integration Feasibility}
\label{sec:integration}
% ─────────────────────────────────────────────────────────────────────────────

\begin{table}[H]
  \centering
  \small
  \caption{Per-engine feasibility matrix for PROBE primitives.}
  \label{tab:engines}
  \begin{tabularx}{\textwidth}{@{}Xccccl@{}}
    \toprule
    Engine & PROBE-RM-final & PROBE-MT & PROBE-HD & B-PROBE-LOGIT & Notes \\
    \midrule
    transformers + bitsandbytes NF4 & \YES & \YES & \YES & \YES & \texttt{register\_forward\_hook}; \textbf{used in this paper} \\
    transformers + autoawq INT4 & \YES & \YES & \YES & \YES & identical hook API \\
    vLLM & \YES & \YES (model-runner patch) & \YES & \YES & logits processor for output \\
    llama.cpp (eval-callback) & \YES & \YES & \YES & \YES & runtime tensor-name matching \\
    llama.cpp (default) & \YES & \NO & \NO & \YES & \texttt{llama\_get\_embeddings\_ith} only \\
    llama-cpp-python (default) & \YES & \NO & \NO & \YES & \texttt{embedding=True} mode \\
    Apple MLX & \YES & \YES & \YES & \YES & module instrumentation \\
    RK3588 RKNN W8A8 & \PART & \NO & \NO & \YES & conversion-time output mark required \\
    Qualcomm QNN / AI Hub & \NO / \PART & \NO & \NO & \YES & frozen graph; recompile for \PART \\
    Google AICore (Gemini Nano) & \NO & \NO & \NO & \YES & closed-graph; B-PROBE only \\
    OpenAI / Anthropic / Gemini API & \NO & \NO & \NO & \YES & top-$k$ logprob API only \\
    \bottomrule
  \end{tabularx}
\end{table}

\noindent\YES{} = feasible; \PART{} = feasible with configuration; \NO{} = infeasible by default.

% ─────────────────────────────────────────────────────────────────────────────
\section{Limitations and Open Questions}
\label{sec:limitations}
% ─────────────────────────────────────────────────────────────────────────────

\subsection{Known Limitations}

\begin{enumerate}[leftmargin=2em, itemsep=4pt]
  \item \textbf{Single model family, two schemes (current).} H1 requires $\ge 2$ of 3
        families; only Llama-3.2-3B at FP16 and NF4 tested.
  \item \textbf{Proxy behavioral metric.} $\bcliff = 0.000$ is computed from
        PROBE-RM margin threshold crossings (partially circular). The full
        StrongREJECT + Llama-Guard-3-8B judge is required for valid H1 acceptance.
  \item \textbf{Cliff zone not tested.} GGUF Q3\_K\_M and Q2\_K (predicted cliff zone
        per~\citet{egashira2024exploiting} and~\citet{mindthegap2025}) have not been
        evaluated. Requires llama.cpp integration.
  \item \textbf{Fold B corpus scope.} HarmBench~\citep{mazeika2024harmbench} and
        ArtPrompt absent; only AdvBench + JBB used.
  \item \textbf{Hardware limitation.} Colab T4 (16\,GB) limits evaluation to 3B models.
        Larger models (8B, 13B) require A100-40G.
  \item \textbf{H2--H5 untested.} Composition (H4), bandit drift (H4/CONDUCTOR),
        Tier-C weakness (H5), black-box FPR decoupling (H3) are all future work.
  \item \textbf{Defender-aware adversaries (A8).} Bandit poisoning is a residual risk.
        Standard mitigation: randomized thresholds + safe rollback. Not formally proven.
  \item \textbf{No certified robustness.} Only FPR-portability theorem and bandit
        regret bounds; no per-input certificate.
  \item \textbf{llama.cpp intermediate layers.} Multi-layer PROBE-MT requires
        eval-callback (PR \#11274) or a fork-and-patch build. Not tested.
  \item \textbf{Calibration drift.} A new firmware shipping a re-quantized model
        invalidates the calibration table. Automatic re-calibration on first boot is
        proposed but not implemented.
\end{enumerate}

\subsection{Open Research Questions}

\begin{itemize}[leftmargin=2em, itemsep=2pt]
  \item Does $\cliff$ show a phase transition or a smooth curve across schemes?
  \item Is there a single $(\rhat, \ell)$ that transfers across quantization schemes
        within one base model (``quant-invariant refusal direction'')?
  \item Can PROBE-RM be augmented with refusal-direction \emph{redundancy}
        (ensembled directions) for cliff resilience?
  \item How tight is the geometric--behavioral correlation at the cliff, and what
        mechanism explains residual disagreement?
  \item Can the bandit's reward signal be made denser via synthetic-injection
        sentinels without storing payloads?
\end{itemize}

% ─────────────────────────────────────────────────────────────────────────────
\section{Evaluation Roadmap}
\label{sec:roadmap}
% ─────────────────────────────────────────────────────────────────────────────

\subsection{Completed (Phase A/B)}

\begin{itemize}[leftmargin=2em]
  \item Full package scaffolding: 8 components, 939 tests, mypy strict.
  \item Fold A + Fold B live run: \texttt{Llama-3.2-3B-Instruct}, FP16 + NF4, Colab T4.
  \item Paper draft (main.tex), figures (6 publication-quality), refs.bib (23 entries).
  \item GitHub Actions workflow for figure regeneration.
\end{itemize}

\subsection{Phase C: GGUF + Multi-Family + Full Judge (1--2 months)}

\begin{enumerate}[leftmargin=2em, itemsep=3pt]
  \item Install llama.cpp with eval-callback on an A100 or RTX GPU.
  \item Evaluate GGUF Q4\_K\_M, Q3\_K\_M (primary H1 target), Q2\_K.
  \item Integrate StrongREJECT + Llama-Guard-3-8B for non-circular $\bcliff$.
  \item Add Qwen-2.5-7B-Instruct and Mistral-7B-Instruct as second and third families.
  \item Run Fold C (defense composition): GCG, AutoDAN, PAIR, TRIPWIRE-H ablation.
  \item Reproduce Egashira-style A7 attack at Q3\_K\_M and test LADDER + ATTEST.
\end{enumerate}

\subsection{Phase D: Bandit + BCN-2 (2--4 months)}

\begin{enumerate}[leftmargin=2em, itemsep=3pt]
  \item Wire CONDUCTOR with real reward signals from LOOKOUT.
  \item Fold D: streaming mixed traffic with adversarial drift episodes (A8).
  \item Build BCN-2 dataset (Mistral-7B-base paraphraser, FP16-refusal filter).
  \item Fold E: BCN-2 cross-family evaluation.
  \item Test H2--H5 with sufficient statistical power.
\end{enumerate}

\subsection{Resource Estimates}

\begin{table}[H]
  \centering
  \small
  \caption{Resource estimates per phase (rough order of magnitude).}
  \label{tab:resources}
  \begin{tabular}{llll}
    \toprule
    Phase & Hardware & Estimated Cost & Time \\
    \midrule
    A/B (done) & Colab T4 (free tier) & $\le \$0$ & 2 sessions \\
    C (GGUF + multi-family) & A100 40G or RTX 4090 & $\le \$100$ cloud & 1--2 weeks \\
    D (bandit + BCN-2) & A100 40G & $\le \$200$ cloud & 2--4 weeks \\
    Tier B reproduction & Raspberry Pi 5 8 GB & $\approx \$120$ hardware & 1 week \\
    Tier C reproduction & RK3588-class SBC & $\approx \$150$ hardware & 1 week \\
    \bottomrule
  \end{tabular}
\end{table}

% ─────────────────────────────────────────────────────────────────────────────
\section{Preregistration and Reproducibility}
\label{sec:preregistration}
% ─────────────────────────────────────────────────────────────────────────────

Before unblinding any results, the following will be published on OSF + arXiv:

\begin{itemize}[leftmargin=2em, itemsep=2pt]
  \item Five-fold split with all SHA-256 hashes.
  \item KenLM training corpora SHA-256.
  \item Judge prompts (StrongREJECT + Llama-Guard-3-8B) verbatim.
  \item Statistical analysis plan including the power calculation (\S\ref{sec:folds}).
  \item All primitive threshold derivation rules.
  \item BCN-2 paraphraser checkpoint hash and generation seeds.
  \item Pre-registered reward function (Eq.~\ref{eq:reward}).
\end{itemize}

\textbf{Code release (phased).}
Phase 1 (paper acceptance): defense code, calibration scripts, integration shims.
Phase 2 (90 days post): Folds A + C + D.
Phase 3 (12 months, gated): BCN-2 (Fold E), hashed-only with controlled institutional access.

% ─────────────────────────────────────────────────────────────────────────────
\section{Responsible Disclosure and Dual-Use}
\label{sec:disclosure}
% ─────────────────────────────────────────────────────────────────────────────

Three dual-use artifacts:

\begin{enumerate}[leftmargin=2em, itemsep=3pt]
  \item \textbf{BCN-2 dataset.} Below-cliff-natural prompts.
        Hashed-only public release; clear-text via gated institutional access.
  \item \textbf{Cliff-localization protocol.} The $\bcliff$ measurement is a
        fingerprint of the cliff.
        Coordinated disclosure with model vendors (Meta, Alibaba, Mistral, Google)
        $\ge 60$ days before public release of any empirically detected cliff cell.
  \item \textbf{Adaptive attacks discovered during evaluation.} Any new attack vectors
        uncovered are reported to relevant vendors and to HarmBench / JBB / AgentDojo
        maintainers before publication.
\end{enumerate}

Disclosure norm: vendors acknowledge within 14 days; mitigations tracked; public
release follows mitigation deployment or 90 days, whichever is sooner.

% ─────────────────────────────────────────────────────────────────────────────
\section{Conclusion}
\label{sec:conclusion}
% ─────────────────────────────────────────────────────────────────────────────

\CG{} reframes prompt-injection defense as a problem in
\textbf{quantization-aware control}: cheap, stateless gates whose thresholds are
pinned to the protected model's own per-quantization residual-stream geometry,
sequenced by a contextual bandit, monitored by streaming change-point detectors,
and ported across hardware tiers via a static router.

As of May 2026, Phase A is complete (939 tests passing, full scaffolding) and
Phase B delivers the first live Fold A/B empirical results:
$\cliff = 0.167$ ($0.67\kap$), $\wcliff = 0.014$, $\bcliff = 0.000$ (proxy)
under NF4 for Llama-3.2-3B-Instruct.
H1 is not accepted for this configuration.
The predicted cliff zone (GGUF Q3\_K\_M and below) is the primary target for
the next evaluation phase.

The theoretical foundation --- FPR-decoupling theorem with black-box corollary,
CUSUM-optimal streaming detection, LinUCB/EXP3.S regret bounds, and the
BCN-2 anti-circularity protocol --- is complete and pre-registered.
What remains is empirical: running Folds C--E at the predicted cliff zone,
with a non-circular behavioral judge, across three model families and five
quantization schemes.

If the cliff does not exist, the experiment falsifies the design center and
we report that.
If it does, \CG{} is the appropriate first defense for the world that
Egashira, Hong, and Mind the Gap have described.

% ─────────────────────────────────────────────────────────────────────────────
\bibliographystyle{plainnat}
\bibliography{refs}

\end{document}
